\documentclass[a4paper]{article}
\usepackage{a4wide}
\usepackage{marvosym}
\usepackage{xspace}
\usepackage{graphics}
\usepackage{color}
\usepackage[colorlinks,linkcolor=blue]{hyperref}
\usepackage[english]{babel}
\pagestyle{empty}
\voffset=-2cm
\addtolength{\textheight}{10cm}

\newcommand{\skype}{\href{http://www.skype.com}{\raisebox{-.11cm}{\resizebox{.3cm}{!}{\includegraphics{skype}}}}\xspace}
\newcommand{\myhref}[2]{\href{#1}{\textcolor{blue}{#2}}}

\begin{document}

\rightline{\parbox{5cm}{\small
\begin{tabular}{ll}
Helsinki, Finland & \Email\myhref{mailto:vitaly.repin@gmail.com}{vitaly.repin@gmail.com}\\
\skype \myhref{skype:vitaly_repin?call}{vitaly\_repin}&\Telefon \myhref{callto://+358504876392}{+358 504876392}
\end{tabular}}}

{\Huge\bf Vitaly Repin}

\newcommand{\tl}[1]{\centerline{\textbf{\textbf{\Large #1}}}}
\vspace*{.8cm}
\begin{tabular}{lp{12cm}}
\bf Objective & Seeking a challenging position as an architect, technical team lead.\\
\bf Experience &\textbf{05.2012 -- Now. CBTec Oy. Helsinki, Finland. Chief software architect, co-founder}\\[-.5cm]
& \begin{itemize}
\item Architecture and project management. CBTec Oy is a new startup company developing product in educational area.
\end{itemize}\\[-.3cm]
&\textbf{2008 -- 2012. Nokia Oyj, Maemo and MeeGo devices. Helsinki, Finland. Software architect, Technical Product Owner}\\[-.5cm]
& \begin{itemize}
\item \textbf{E-Mail and MfE (Exchange ActiveSync) subsystems for \myhref{http://en.wikipedia.org/wiki/Nokia_N9}{Nokia N9.}} Qt-based MeeGo/Harmattan E-Mail client.
Fully responsible for the software architecture and project execution.
\item \textbf{E-Mail subsystem for \myhref{http://en.wikipedia.org/wiki/N900}{Nokia N900.}} Architecture routines and project management. Open source e-mail client \myhref{http://modest.garage.maemo.org/}{modest.}
\item \textbf{MailForExchange subsystem for \myhref{http://en.wikipedia.org/wiki/N900}{Nokia N900.}} Bringing MfE support on board in 4 months with 0 crashes in sales
release. Defining the architecture, leading the team. Cooperation with multiple geographically-spreaded teams.
After sales release in addition to my engineering responsibilities I was communicating with customers through \myhref{http://talk.maemo.org/}{maemo.org} (public channel) and Nokia Care (private channel).
\item \textbf{QMF (Qt Messaging Framework)} Bringing on board in Nokia Linux-based platform "Harmattan" MeeGo. Deployment (storage part) in Nokia N900 device (first Linux phone made by Nokia).
\item \textbf{Processes:} Architecture review process, established code review process for email team, hiring interviews with the candidates for technical positions.
\item \textbf{Technical contribution:} participating in Harmattan MeeGo architecture board (influencing the Harmattan
MeeGo architecture, leading crash meetings (core dump analysis), taking the lead in analysis of the most complex
issues with the email and mfe subsystem (performance, battery life, hangs, rare crashes etc), defining architecture for the email subsystem in
general and email client in particular, defining architecture for ActiveSync solution.
\end{itemize}\\[-.25cm]
& \textbf{2007 -- 2008. Nokia Oyj, OSSO and Maemo. Helsinki, Finland. Software developer.}\\
& My projects: Skype for N800 and N810. Rhapsody for N800. The roles: software development (C and C++, gtk+/glib),
software design, code reviews.  Cooperation with offshore teams. Troubleshooting (voice quality, backend bugs).
I have also made a proof of concept for bringing SWT with browser integration to the platform (on top of
\myhref{https://evolvis.org/plugins/mediawiki/wiki/jalimo/index.php/Jalimo}{Jalimo}) which was demonstrated by IBM and Nokia
in one of the technological shows in US.\\
& \textbf{2005 -- 2007. Ice Brains Software Co., LLC. St. Petersburg, Russia.  Team leader.}\\
& I acted as team leader for several (software outcourcing) projects related to Skype, Asterisk and \myhref{http://www.ltsp.org/}{LTSP.} Won he Skype prize as "developer of the month" for cross-platform C++ library to work with Skype API.\\
& \textbf{2003 -- 2005. Inotech Co., LLC. St. Petersburg, Russia. Team leader.}\\
& Design and development of the billing system (BSS) PLATEX v.3.  This system is used by the variety of telecommunication operators in Russia and CIS.\\
& \textbf{2001 -- 2003. Inotech Co., LLC. St. Petersburg, Russia.  Senior software developer.}\\
& Design and development of the call center for Metrocom JSK, St.-Petersburg.  
The hardware base for this project was EXS switch.  This switch is connected to the PBX (Siemens EWSD) via SS7 and can be controlled
by the software via TCP/IP. My parts of the project: overall system design, EXS setup and configuring, EXS management via TCP/IP, IVR imlementation.\\
\end{tabular}

\newpage

\begin{tabular}{lp{12cm}}
& \bf 1999 -- 2001. Komin Co., Ltd (St. Petersburg center of satellite communications). St. Petersburg, Russia. Software developer.\\[-.5cm]
& \begin{itemize}
\item Design and implementation
of the remote control system for complex scientific device 
(high-precision spectroanalyzer located near the radiotelescope).
Developed software consists of two parts --- server (implemented as several
UNIX applications) and client (implemented as windows application).
Server and client are connected via the Internet (TCP/IP). 
This project won the grant at the St.-Petersburg grant
competition for students, PhDs, young scientists  and specialists.
\item Design and implementation of the host software for the  
voice compression system prototype based on voice compression codec produced by
\myhref{http://www.audiocodes.com}{AudioCodes Ltd.}
This work was part of IP--telephony project which was developing at 
Komin Ltd.  Besides, it was my bachelor work.  
This work got diploma and RosTelecom company premium at the All-Russian 
competition of student's researches.
Besides, this work won the grant at the St.-Petersburg grant competition
for students, PhDs, young scientists  and specialists.
\item Design and implementation of the software for data aquisition board
developed at Komin Ltd.
Developed software consists of two parts --- driver (implemented as loadable 
kernel module for RT Linux) and user--level application (GSM (voice and signaling) decoder with
Qt-based GUI).
\item Design and implementation of software DTMF 
detector.  The result of this work was \texttt{TDTMFDetector}
class which provided the possibility
to detect any in--band signal encoded by frequency pair $(f_1, f_2)$ and
to measure its duration.
\end{itemize}\\
\bf Education& \bf 2002  Radiophysical Science \& Engineering (RFS\&E)\\
& \bf Faculty of the  St. Petersburg State Polytechnical\\
& \bf University, Russia.\\
& Master of Science.\\
\bf Academic\\
\bf Honor & \vspace*{-1cm}\begin{itemize}
\item Two grants at the St.-Petersburg grant competition
for students, PhDs, young scientists  and specialists.
\item My bachelor work got diploma at the All-Russian competition of student's researches.
\end{itemize}
\end{tabular}

\begin{tabular}{lp{12cm}}
\bf Skills& \vspace*{-.6cm}\begin{itemize}
\item \textbf{Software process:} Extreme Programming (XP), SCRUM
\item \textbf{Programming languages:} C, C++, Pascal, Tcl/Tk, Perl, assemblers (x86, ADSP).
\item \textbf{OS:}  Linux, RTLinux, Windows, DOS.
\item \textbf{Databases:} Oracle, MySQL.
\item \textbf{Applications:} \LaTeXe, Asterisk, LTSP (Linux Terminal Server Project), Apache, Samba, MTAs...  
\item \textbf{Other:}     OOA/OOD/OOP, SQL, shell scripting languages, 
			theory and practice of DSP, VoIP (SIP, MGCP, H.323, Skype), 
			basis of SS7, very expirienced with the Excel Switching System (EXS),
			OSS/BSS development expirience.
\end{itemize}\\
\bf Personal\\
\bf data& \vspace*{-1cm}\begin{itemize}
\item \textbf{Birth date:} August 18, 1979
\item \textbf{Marital status:} Single
\item \textbf{Nationality:} Russian
\item \textbf{Resident of:} Finland
\item \textbf{Hobby:}  Outdoor activities, bicycle
\end{itemize}\\
\bf Personal\\
\bf qualities & \vspace*{-.7cm}
Creative, executive, ability to learn new technologies/projects quickly,
ability to work under pressure, experienced team lead.
\end{tabular}

\end{document} 
